%Phần thiết đặt trang
\documentclass[14pt,a4paper]{article}
\usepackage{mathastext}
\usepackage[utf8]{vietnam}
\usepackage{amsfonts}
\usepackage{amsmath}
\usepackage{amssymb}
\usepackage{graphicx}
\usepackage[left=2cm,right=2cm,top=2cm,bottom=2cm]{geometry}
%Thụt đầu dòng cho tất cả
\usepackage{indentfirst}
\setlength{\parindent}{1.5cm}

%\usepackage{parskip}
%\setlength{\parskip}{0.5em}

%Đánh dấu đầu dòng bằng chữ a) b) c)
\usepackage[shortlabels]{enumitem}

%Font chữ hoa văn cho các chữ cái viết hoa
\usepackage[mathscr]{euscript}

%Điều chỉnh kích cỡ chữ cho tài liệu
\usepackage{scrextend}
\changefontsizes[20pt]{14pt}

%Thêm phần header cho trang
\usepackage{fancyhdr}
\pagestyle{fancy}
\renewcommand{\subsectionmark}[1]{\markboth{#1}{}}
\fancyhf{}
\fancyhead[R]{\leftmark}
\fancyfoot[C]{Trang \thepage}
\fancypagestyle{plain}{%
    \fancyhf{}%
    \renewcommand{\headrulewidth}{0pt}%
}

%Thêm đường line dưới footer cho trang
\renewcommand{\footrulewidth}{0.4pt}

%Điều chỉnh nội dung bảng danh sách hình vẽ
\usepackage[titles]{tocloft}
\newlength{\mylen}
\renewcommand{\cftfigpresnum}{\figurename\enspace}
\renewcommand{\cftfigaftersnum}{:}
\settowidth{\mylen}{\cftfigpresnum\cftfigaftersnum}
\addtolength{\cftfignumwidth}{\mylen}


%Màu sắc cho bảng
\usepackage[table,xcdraw]{xcolor}

%Chuyển hình EPS sang dạng PDF tránh gây lỗi font
\usepackage{epstopdf}

%Cải thiện sắp xếp vị trí hình ảnh, dùng kèm với [H] trong phần option ảnh
\usepackage{float}

%Dùng dấu chấm thay dấu hai chấm trong caption ảnh
\usepackage[labelsep=period]{caption}

%Phần thiết kế khung code nhập liệu
\usepackage{listings}
\usepackage{color}

\definecolor{dkgreen}{rgb}{0,0.6,0}
\definecolor{gray}{rgb}{0.5,0.5,0.5}
\definecolor{mauve}{rgb}{0.58,0,0.82}

\lstset{frame=tb,
  language=Matlab,
  aboveskip=3mm,
  belowskip=3mm,
  showstringspaces=false,
  columns=flexible,
  basicstyle={\small\ttfamily},
  numbers=left,
  numberstyle=\small\color{gray},
  keywordstyle=\color{blue},
  commentstyle=\color{dkgreen},
  stringstyle=\color{mauve},
  breaklines=true,
  breakatwhitespace=true,
  tabsize=3
}

%Mở tính năng đánh dấu công thức trong đoạn văn
\makeatletter
\newcommand*{\inlineequation}[2][]{%
  \begingroup
    % Put \refstepcounter at the beginning, because
    % package `hyperref' sets the anchor here.
    \refstepcounter{equation}%
    \ifx\\#1\\%
    \else
      \label{#1}%
    \fi
    % prevent line breaks inside equation
    \relpenalty=10000 %
    \binoppenalty=10000 %
    \ensuremath{%
      % \displaystyle % larger fractions, ...
      #2%
    }%
    ~\@eqnnum
  \endgroup
}
\makeatother

\usepackage{hyperref}

%Phần chọn font chèn câu lệnh giữa đoạn
\newenvironment{code}{\ttfamily}{\par}
\DeclareTextFontCommand{\chuyencode}{\code}

%Thêm chấm vào tiêu đề các phần
\usepackage{titlesec}
\titlelabel{\thetitle.\quad}

%Thêm định dạng cho các nút và menu lệnh.
\usepackage{menukeys}

%Đánh số cho các ví dụ và bài tập
\usepackage[thref,thmmarks,standard,amsmath,hyperref]{ntheorem}
\theoremheaderfont{\bfseries}
\theorembodyfont{\normalfont}
\theoremseparator{:}
\renewtheorem{example}{Ví dụ}
%Đánh số công thức theo chương
\numberwithin{equation}{section}

%Tự ngắt dòng cho biểu thức toán dài dòng
\usepackage{mathpazo}
\usepackage[mathpazo]{flexisym}
\usepackage{breqn}

%Khung bìa
\usepackage{fancybox}

%Nội dung chính
\begin{document}
%Chỉnh loại tiêu đề chương thành I, II
\renewcommand\thesection{\Roman{section}}
\renewcommand\thesubsection{\arabic{subsection}}
\renewcommand\thesubsubsection{\thesubsection.\arabic{subsubsection}}
%Trang bìa
\thispagestyle{empty}
\thisfancypage{
\setlength{\fboxsep}{2pt}
\fbox}{} 
\begin{center}
\begin{large}
\textbf{\\TRƯỜNG ĐẠI HỌC CẦN THƠ}
\end{large} \\
\begin{large}
\textbf{KHOA KHOA HỌC TỰ NHIÊN}
\end{large} \\
\textbf{--------------------  *  ---------------------}\\[1cm]
\includegraphics{logo-ctu} \\[2cm]
{\fontsize{18pt}{1}\selectfont TIỂU LUẬN\\HỌC PHẦN "QUÁ TRÌNH NGẪU NHIÊN"}\\[2cm]
{\fontsize{22pt}{1}\selectfont \textbf{PHƯƠNG TRÌNH VI PHÂN NGẪU NHIÊN}}\\[3cm]
\end{center}


\indent \begin{tabular}{ll}
Giáo viên hướng dẫn: &\textbf{TS. Dương Tôn Đảm}\\
Học viên thực hiện : &\textbf{Nguyễn Dương Thanh Phú} \\
MSHV: &\textbf{M1817008}\\
Lớp: &\textbf{LTXS và Thống kê Toán học K24}
\end{tabular}


\vspace{4cm}
\begin{center}
{\fontsize{16pt}{1}\selectfont CẦN THƠ}\\
{\fontsize{16pt}{1}\selectfont \today}
\end{center}
\newpage
%Nội dung chính
\section*{\begin{center}
	LỜI NÓI ĐẦU
\end{center}}
\addcontentsline{toc}{section}{Lời nói đầu}
Phương trình vi phân ngẫu nhiên đóng vai trò rất quan trọng trong kỹ thuật, vật lý, kinh tế và một số ngành khoa học khác. Sự ra đời của nó xuất phát từ hu cầu xác dịnh mối quan hệ giữa một bên là đại lượng biến thiên liên tục với một bên là độ biến thiên của đại lượng đó. Các mối quan hệ như thế xuất hiện thường xuyên trong các ứng dụng thực thế. Chẳng hạn, các biến động giá cả của một loại cổ phiếu trên thị trường hay sự thay đổi biến thiên của triều cường trong một khu vực bến cảng... Phương trình vi phân ngẫu nhiên luôn có mối quan hệ chặt chẽ với phép tính toán ngẫu nhiên và tích phân ngẫu nhiên để hình thành nên cơ sở của giải tích ngẫu nhiên. Nó còn hiện diện trong những chủ đề thời sự của giải tích, xác suất và được sử dụng rộng rãi trong nhiều lĩnh vực bên ngoài Toán học.

Thông qua bài tiểu luận này, tôi xin trình bày những nội dung cơ bản nhất về phương trình vi phân ngẫu nhiên và những dạng phương trình đặc trưng phổ biến nhất. Đi kèm với những định lý và chứng minh quan trọng làm nền tảng cho mọi người để bắt đầu tìm hiểu về phương trình vi phân ngẫu nhiên.

Mặc dù đã cố gắng hết sức để hoàn thành này, nhưng tiểu luận sẽ không thể tránh khỏi những thiếu sót, rất mong thầy cô và bạn đọc đóng góp và bổ sung ý kiến để tiểu luận có thể hoàn thiện hơn nữa.

Chân thành cảm ơn TS. Dương Tôn Đảm đã hết lòng hỗ trợ về đề tài lẫn nội dung để tôi hoàn thành tiểu luận. Đồng thời, cảm ơn bạn đọc đã dành thời gian đọc và góp ý cho tiểu luận này.
\newpage
%Chương 1:
\section{Các kiến thức chuẩn bị}
\subsection{Mở đầu và ví dụ về phương trình vi ngẫu nhiên}
\subsubsection{Mở đầu}
Trước tiên, ta tìm các quá trình $X=(X_t,t\geq 0)$ có thể thoả mãn phương trình
\begin{equation}\label{eq:1.1}
	\dfrac{dX_t}{dt}=b(t,X_t)+\sigma(t,X_t)W_t
\end{equation}
trong đó $b(t,x)\in\mathbb{R},\sigma(t,x)\in\mathbb{R}$ và $W_t$ là quá trình Wiener 1-chiều. Khi đó $X_t$ cũng là lời giải của phương trình tích phân
\begin{equation}\label{eq:1.2}
	X_t=X_0+\int_0^tb(s,X_s)ds+\int_0^t\sigma(s,X_s)dB_s
\end{equation}
hay viết dưới dạng vi phân là
\begin{equation}\label{eq:1.3}
	dX_t=b(t,X_t)dt+\sigma(t,X_t)dB_t
\end{equation}
Như vậy từ \eqref{eq:1.1} sang \eqref{eq:1.3} ta đã thay thế một cách hình thức $W_t$ bởi $\dfrac{dB_t}{dt}$. Vấn đề đặt ra lúc này với một phương trình vi phân \eqref{eq:1.3} thì:
\begin{enumerate}[(1)]
	\item Với những điều kiện nào của các hệ số b và $\sigma$ thì tồn tại nghiệm và nghiệm đó là duy nhất?
	\item Ta sẽ giải phương trình đó như thế nào?
\end{enumerate}
Ta hãy xét câu hỏi 2 qua một số ví dụ khác, sau đó ở phần sau sẽ giải quyết vấn đề trong câu hỏi 1.
\subsubsection{Một số ví dụ về phương trình vi phân ngẫu nhiên}
\textbf{a. Mô hình tăng dân số đơn giản}
\begin{equation}\label{eq:1.1.a}
	\dfrac{dN_t}{dt}=\alpha_t.N_t \condition{$N(0)=N_0$}
\end{equation}
trong đó $N_t$ là số dân tại thời điểm $t,\alpha_t$ là tốc độ tăng dân số tại thời điểm t. Thông thường $\alpha_t$ không được xác định rõ ràng nhưng do chịu tác động của môi trường tự nhiên nên người ta thường viết dưới dạng $\alpha W_t$, với $W_t$ là một quá trình Wiener còn $\alpha$ là hằng số. Ngoài ra, khi kết hợp với tỉ lệ sinh và tử cho ta thêm một hằng số khác là $r_t$. Do đó ta được $\alpha_t=r_t+\alpha W_t$ và thay vào phương trình \eqref{eq:1.1.a} sẽ được:
\begin{equation*}
	dN_t=rN_tdt+\alpha N_tdB_t
\end{equation*}
trong đó cụm $b(t,x)=rx$, còn $\sigma(t,x)=\alpha x$ hay
\begin{equation*}
	\dfrac{dN_t}{N_t}=rdt+\alpha dB_t\Leftrightarrow\int_0^t\dfrac{dN_t}{N_t}=rt+\alpha B_t \condition{$B_0=0$}\tag{*}
\end{equation*}
Để tính tích phân vế trái, ta áp dụng công thức Itô cho hàm $g(t,x)=\ln(x),x>0$ và ta có:
\begin{dmath*}
	d(\ln N_t)=\dfrac{1}{N_t}dN_t+\dfrac{1}{2}(-\dfrac{1}{N_t^2})(dN_t)^2=\dfrac{dN_t}{N_t}-\dfrac{1}{2N_t^2}\alpha^2N_t^2dt=\dfrac{dN_t}{N_t}-\dfrac{1}{2}\alpha^2.dt
\end{dmath*}
Do đó
\begin{equation*}
	\dfrac{dN_t}{N_t}=d(\ln N_t)+\dfrac{1}{2}\alpha^2dt
\end{equation*}
Kết hợp với hệ thức (*) ta suy ra:
\begin{equation*}
	\ln\dfrac{N_t}{N_0}=(r-\dfrac{1}{2}\alpha^2)\alpha t+\alpha B_t
\end{equation*}
Hay:
\begin{equation*}
	N_t=N_0.exp[(r-\dfrac{\alpha^2}{2})t+\alpha B_t]
\end{equation*}
Đó là lời giải của mô hình dân số đã cho.\\
\textbf{b. Chuyển động Brown trên đường tròn đơn vị}\\
Ta chọn $X=B$, một chuyển động Brown 1 chiều và $g(t,x)=e^{ix}=(\cos{x},\sin{x})\in\mathbb{R}^2$ với $x\in\mathbb{R}$. Khi đó
\begin{equation*}
	Y\equiv g(t,X)=e^{iB}=(\cos{B},\sin{B})
\end{equation*} 
Vậy
\begin{equation*}
	Y=(Y_1,Y_2) \condition{với $Y_1=\cos{B},Y_2=\sin{B}$}
\end{equation*}
Theo công thức Itô:
\begin{equation*}
	dY_1=-\sin{B.dB}-\dfrac{1}{2}\cos{B}dt
\end{equation*}
\begin{equation*}
	dY_2=-\cos{B.dB}-\dfrac{1}{2}\sin{B}dt
\end{equation*}
Như vậy, quá trình 2 chiều $Y=(Y_1,Y_2)$ mà ta gọi là chuyển động Brown trên đường tròn đơn vị, $(Y_1^2+Y_2^2=1)$ là lời giải của hệ phương trình vi ngẫu nhiên
\begin{equation*}
	dY_1=-\dfrac{1}{2}Y_1dt-Y_2dB,
\end{equation*}
\begin{equation*}
	dY_2=-\dfrac{1}{2}Y_2dt+Y_1dB
\end{equation*}
hoặc dưới dạng ma trận
\begin{equation*}
	dY=-\dfrac{1}{2}Y+KYdB,
\end{equation*} 
trong đó
\begin{equation*}
	Y=(Y_1,Y_2)\condition{$K=\begin{pmatrix} 0 & -1 \\ 1 & 0 \end{pmatrix}$}.
\end{equation*}
\subsection{Điều kiện về sự tồn tại và duy nhất của nghiệm}
\subsubsection{Định lý về sự tồn tại và duy nhất lời giải}
\textbf{a. Phát biểu định lý}\\
Cho
\begin{align}
	& \begin{aligned}
	T>0,& b:[0,T]\times\mathbb{R}^n\rightarrow\mathbb{R}^n\\
	& \sigma:[0,T]\times\mathbb{R}^n\rightarrow\mathbb{R}^{n\times m}
\end{aligned}
\end{align}
là những hàm đo được và thoả mãn
\begin{enumerate}[(i)]
	\item $|b(t,x)|+|\sigma(t,x)\leq C(1+|x|)|$, với $x\in\mathbb{R}^n,t\in[0,T]$ và với hằng số C nào đó, trong đó $|b|$ là độ dài vectơ còn $|\sigma|=\sum{|\sigma_{ij}|^2}$.
	\item $|b(t,x)-b(t,y)|+|\sigma(t,x)-\sigma(t,y)|\leq D.|x-y|,x,y\in\mathbb{R}^n,t\in[0,T]$ và D là hằng số nào đó.
\end{enumerate}
Cho $Z$ là một biến ngẫu nhiên độc lập với $\sigma$-trường $\mathscr{F}_\infty=\sigma \condition{$(B_t;0<t<\infty)$}$ sao cho $E(|Z|^2)<\infty$. Kh đó, phương trình vi phân ngẫu nhiên
\begin{equation*}
	dX_t=b(t,X_t)dt+\sigma(t,X_t)dB_t \tag{PT}
\end{equation*}
\begin{equation*}
	X_0=Z
\end{equation*}
có một lời giải duy nhất, liên tục theo t, mỗi thành phần của $X_t$ là một quá trình thuộc về lớp $N[0,T]=\{f:[0,\infty)\times\Omega\rightarrow\mathbb{R},\mathscr{B}\times\mathscr{F}$-đo được,$\mathscr{F}_t$-thích nghi và $E\int_0^T|f(t,\omega)|^2dt<\infty$, trong đó $\mathscr{B}=\mathscr{B}_{[0,\infty)},\mathscr{F}_t=\sigma(B_s,s\leq t)\}$\\
\indent\textbf{b. Giải thích}\\
Ta sẽ giải thích hai điều kiện (i) và (ii) trong phát biểu định lý nói trên qua hai ví dụ sau đây và sẽ dễ dàng thấy rằng đó đều là những điều kiện rất tự nhiên. Các ví dụ này sẽ liên quan đến các phương trình vi phân tất định (khi đó $\sigma=0$).\\
\begin{example}
Xét phương trình $\dfrac{dX_t}{dt}=X^2_t,X_0=1$ (tức là $dX_t=X_t^2dt,X_0=1$).	
\end{example}
Ở đây $b(s)=x^2$ không thoả mãn điều kiện (i). Phương trình này có một nghiệm duy nhất
\begin{equation*}
	X_t=\dfrac{1}{1-t} \condition{$t\neq 1$}
\end{equation*}
Nghiệm này không xác định với $t=1$ và $X_t\rightarrow\infty$ khi $t\rightarrow 1$ (ta thường nói: $X_t$ "bùng nổ" hay xuất phát tại $t=1$).\\
Nói chung, điều kiện (i) dùng để đảm bảo cho lời giải của (PT) $X_t(\omega)$ không bùng nổ, tức là $|X_t(\omega)|\not\to\infty$ tại một điểm hữu hạn.\\
\begin{example}
	Xét phương trình $\dfrac{dX_t}{dt}=3X^{2/3}_t,X(0)=0$ (tức $dX_t=3X^{2/3}_tdt,X_0=1$).
\end{example}
Ở đây $b=3x^{2/3}$ không thoả mãn điều kiện Lipschitz (ii). Từ đó ta có thể suy ra một lời giải của phương trình trên là $X_t=t^3$. Nhưng phương trình nào còn có nhiều lời giải khác. Xét với mọi $a>0$ ta còn các nghiệm
\begin{equation*}
	X_t=\begin{cases}
		0&\condition{với $t\leq a$}\\
		(t-a)^3&\condition{với $t>a$}
	\end{cases}
\end{equation*}
Vậy điều kiện (ii) đảm bảo tính duy nhất của lời giải của phương trình (PT), theo nghĩa hầu chắc chắn.\\
\indent\textbf{c. Chứng minh Định lý tồn tại duy nhất lời giải}\\
\textbf{$\bullet$ Sự duy nhất}\\
Sự duy nhất sẽ được chứng minh dựa vào sự đẳng cự Itô và Điều kiện Lipschitz (ii). Giả sử
\begin{equation*}
	X_1(t,\omega)=X_t(\omega)
\end{equation*}
và
\begin{equation*}
	X_2(t,\omega)=\widehat{X}_t(\omega)
\end{equation*}
là hai lời giải với các điều kiện ban đầu $Z$ và $\widehat{Z}$, tức là $X_1(0,\omega)=Z(\omega),X_2(0,\omega)=\widehat{Z}(\omega),\omega\in\Omega$.\\
Thực ra ở đây ta chỉ cần tới trường hợp $Z=\widehat{Z}$ nhưng ta sẽ đưa vào một ước lượng tổng quát hơn sau đây, liên quan đến tính liên tục Feller.\\
Đặt $a(s,\omega)=b(s,X_s)-b(s,\widehat{X}_s)$ và $\gamma(s,\omega)=\sigma(s,X_s)-\sigma(s,\widehat{X}_s)$. Khi đó ta có:
\begin{equation*}
\begin{split}
	E[|X_t-\widehat{X}_t|^2]&=E[(Z-\widehat{Z}+\int_0^tads+\int_0^t\gamma dB_s)^2]\\
&\leq 3E[|Z-\widehat{Z}|^2]+3E|(\int_0^tads)^2]+3E|(\int_0^t\gamma dB_s)^2]\\
&\leq 3E|Z-\widehat{Z}|^2+E(1+t).D^2\int_0^tE|X_s-\widehat{X}_s|^2ds
\end{split} \condition{(BĐT Cauchy)}
\end{equation*}
bởi vì: theo (ii) $|a|+|\gamma|\leq D.|X-\widehat{X}|$, do đó: $a^2$ hoặc $\gamma^2\leq a^2+\gamma^2\leq (|a|+|\gamma|)^2\leq D^.|X-\widehat{X}_s|^2)$. Vậy hàm
\begin{equation*}
	v(t)=E|X_t-\widehat{X}_t|^2\condition{$0\leq t\leq T$}
\end{equation*}
thoả mãn
\begin{equation*}
	v(t)\leq F+A.\int_0^tv(s)ds,
\end{equation*}
trong đó 
\begin{equation*}
	F=3E|Z-\widehat{Z}|^2\text{ và }A=3(1+T)D^2
\end{equation*}
Đặt $\omega(t)=\int_0^tv(s)$. Khi đó $\omega^'(t)\leq F+A.\omega(t)$. Do đó, vì $\omega(0)=0$ nên $\omega(t)\leq\dfrac{F}{A}(e^{At}-1)$. Theo bất đẳng thức tìm được ở trên $v(t)\leq F+A\omega(t)$, ta có
\begin{equation*}
	v(t)\leq F.exp(At)
\end{equation*}
hay là
\begin{equation*}
	E|X_t-\widehat{X}_t|^2\leq F.exp(At)
\end{equation*}
Bây giờ ta giải thiết rằng $Z=\widehat{Z}$. Khi đó $F=0$ và do đó $v(t)=0$ với mọi $t\geq 0$. Từ đó ta có
\begin{equation*}
	P[|X_t-\widehat{X}_t|=0,\forall t\in \mathbb{Q}\cap[0,T]]=1
\end{equation*}
trong đó $\mathbb{Q}$ là tập hợp các số hữu tỉ. Do tính liên tục hầu khắp nơi của ánh xạ $t\rightarrow|X_t-\widehat{X}_t|$ ta suy ra rằng
\begin{equation*}
	P\{X_t(\omega)=\widehat{X}_t(\omega),\forall t\in[0,T]\}=1
\end{equation*}
hay
\begin{equation*}
	X_1(t,\omega)=X_2(t,\omega)\condition{h.c.c}
\end{equation*}
\textbf{$\bullet$ Sự tồn tại}\\
Ta sẽ chứng minh sự tồn tại lời giải của (PT) theo một phương pháp tương tự đối với phương trình vi phân thường, bằng cách dùng phép lặp Picard. Trước tiên, ta định nghĩa: $Y_t^{(0)}=X_0$ và $Y_t^{(k)}=Y_t^{(k)}(\omega)$ một cách quy nạp như sau:
\begin{equation}\label{eq:stt1}
	Y_t^{(k+1)}=X-0+\int_0^tb(s,Y_s^{(k)})ds+\int_0^t\sigma(s,Y_s^{(k)})dB_s
\end{equation}
Khi đó, tính toán tương tự như đối với Phần Duy nhất ở trên, ta có:
\begin{equation}\label{eq:stt2}
	E|Y_t^{(k+1)}-Y_t^{(k)}|\leq(1+T)3D^2\int_0^tE|Y^{(k)}_s-Y^{(k-1)}_s|^2ds
\end{equation}
với $k\geq 1,t\leq T$ và
\begin{equation*}
	\begin{split}
		E|Y^{(1)}_t-Y^{(0)}_t|^2&\leq 2C^2t^2(1+EX_0^2)+2C^2t(1+EX^2_0)\\
		&\leq 2C^2T.t(1+EX^2_0)+2C^2t(1+EX_0^2)=A_1t,
	\end{split}
\end{equation*}
trong đó hằng số $A_1$ chỉ phụ thuộc vào $C,T$ và $EX^2_0$. Do đó, theo \eqref{eq:stt2} ta có:
\begin{equation*}
	E|Y_t^{(2)}-Y_t^{(1)}|^2\leq(1+T).3D^2.\int_0^tA_1.tdt\underbrace{(1+T).3D^2.A_1}_{A_2}.\dfrac{t^2}{2!}
\end{equation*}
Quy nạp theo k, ta có:
\begin{equation}\label{eq:stt3}
	E|Y_t^{(k+1)}-Y_t^{(k)}|^2\leq\dfrac{A_2^{(k+1)}.t^{(k+1)}}{(k+1)!}\condition{$k\geq 0,t\in[0,T]$}
\end{equation}
trong đó $A_2$ là một hằng số chỉ phụ thuộc $C,D,T$ và $EX_0^2$.\\
Tiếp theo, với mỗi $\omega$ cố định thuộc $\Omega$, ta có:
\begin{equation*}
	\begin{split}
		\sup_{0\leq t\leq T}|Y_t^{(k+1)}-Y_t^{(k)}|&\leq\int_0^t|b(s,Y_s^{(k)})-b(s,Y_s^{(k-1)})|ds+\\
		&+\sup_{0\leq t\leq T}|\int_0^t[\sigma(s,Y_t^{(k)})-\sigma(s,Y_s^{(k-1)})]dB_s|
	\end{split}
\end{equation*}
Theo bất đẳng thức Martingale của Doob
\begin{equation*}
	P\left(\sup_{0\leq t\leq T}|M_t|\geq\lambda \right)\leq\dfrac{1}{\lambda^p}E|M_t|^p
\end{equation*}
thì ta có
\begin{equation*}
\begin{split}
P\left(\sup_{0\leq t\leq T}|Y_t^{(k+1)}-Y_t^{(k)}|\geq 2^{-k} \right)&\leq P\left\{\left[\int\limits_0^T(b(s,Y_s^{(k)})-b(s,Y_s^{(k-1)}))ds \right]^2>\left(\dfrac{1}{2^{k+1}}\right)^2 \right\}\\
&+P\{\sup{|\underbrace{\int\limits_0^t(\sigma(s,Y_s^{(k)})-\sigma(s,Y_s^{(k-1)}dB_s}_{mactingan}|}>\dfrac{1}{2^{(k+1)}} \}\\
&\leq	2^{2k+2}.T.\int_0^TE|b(s,Y_s^{(k)})-b(s,Y_s^{(k-1)})|^2ds\\
&+2^{2k+2}.\int_0^TE|\sigma(s,Y_s^{(k)})-\sigma(s,Y_s^{(k-1)})|^2ds\\
&\leq 2^{2k+2}.D^2(T+1).\int_0^T\dfrac{A_2^kt^k}{k!}dt\\
&\leq\dfrac{(4A_2T)^{k+1}}{(k+1)!}\text{ nếu }A_2\geq4D^2(T+1)
\end{split}
\end{equation*}
Do đó bổ đề Borel-Cantelli, ta có
\begin{equation*}
	P(\sup{|Y_t^{(k+1)}-Y_t^{(k)}|}>\dfrac{1}{2^k} \text{ với vô số k)}=0
\end{equation*}
Vậy với hầu hết $\omega$, tồn tại một $k_0=k_0(\omega)$ sao cho
\begin{equation*}
	\sup_{0\leq t\leq T}|Y_t^{(k+1)}-Y_t^{(k)}|<\dfrac{1}{2^k}\text{ với }k>k_0(\omega)
\end{equation*}
Do đó dãy
\begin{equation*}
	Y_t^{(n)}(\omega)=Y_t^{(0)}(\omega)+\sum_{k=0}^{n-1}[Y_t^{(k+1)}(\omega)-Y_t^{(k)}(\omega)]
\end{equation*}
hội tụ đều với hầu hết mọi $\omega$ (hội tụ điều hầu khắp nơi). Ký hiệu giới hạn đó bởi $X_t=X_t(\omega)$. Khi đó:
\begin{equation*}
	X_t\begin{cases}
		\text{là $t$-liên tục với hầu hết $\omega$ bởi vì $Y_t^{(n)}$ là $t$-liên tục h.k.n với $\forall n$}\\
		\text{là $\mathscr{F}_t$-đo được $\forall t$, vì $Y_t^{(n)}$ là $\mathscr{F}_t$-đo được $\forall t,\forall n$}
	\end{cases}
\end{equation*}
Tiếp theo phần chứng minh là giới hạn trong không gian đủ $L^2(P)$:\\
Cho $m>n\geq0$ là các số nguyên. Theo \eqref{eq:stt3} ta có
\begin{equation*}
	\begin{split}
		E(|Y_t^{(m)}-Y_t^{(n)}|^2)&=||Y_t^{(m)}-Y_t^{(n)}||^2_{L^2(P)}\\
		&=\left[||\sum_{k=n}^{m-1}[Y_t^{(k+1)}-Y_t^{(k)}]||_{L^2(P)}\right]^2\\
		&\leq\left[\sum_{k=n}^{m-1}||Y_t^{(k+1)}-Y_t^{(k)}||_{L^2(P)}\right]^2\\
		&=\left[\sum_{k=n}^{m-1}E(|Y_t^{(k+1)}-Y_t^{(k)}|^2)^{\frac{1}{2}}\right]^2\\
		&\leq\left(\sum_{k=n}^{\infty}\dfrac{(A_2T)^{k+1}}{(k+1)!}\right)\longrightarrow0\text{ khi }n\rightarrow\infty.
	\end{split}
\end{equation*}
Do đó $\{Y_t^{(n)}\}$ hội tụ trong $L_2(P)$ đến một mốc giới hạn mà ta ký hiệu là $Y_t:Y_t^{(n)}\xrightarrow [ (n\rightarrow\infty) ]{ L^2(P) }Y_t$. Vậy tồn tại một dãy con của $\{Y_t^{(n)}(\omega)\}$ hội tụ theo từng điểm đến $Y_t(\omega)$, tức là $|Y_t^{(n)}(\omega)-Y_t(\omega)|\rightarrow0$, và do đó ta phải có $Y_t=X_t$ hầu chắc chắn. Việc còn lại là phải chứng minh rằng $Y_t$ thoả mãn phương trình vi phân (PT) đã cho:\\
Với mọi $n$, ta có:
\begin{equation}\label{eq:stt5}
	Y_t^{(n+1)}=X_0+\int_0^tb(s,Y_s^{(n)})ds+\int_0^t\sigma(s,Y_s^{(n)})dB_s
\end{equation}
Ta đã có $Y_t^{(n+1)}\rightarrow X_t$ khi $n\rightarrow\infty$ và sự hội tụ là điều theo $t\in[0,T]$ với hầu hết $\omega$. Theo (4) và Bổ đề Fatou ta có
\begin{equation*}
	E[\int_0^t|X_t-Y_t^{(n)}|^2dt]\leq\limsup_{m\rightarrow\infty}E\int_0^t|Y_t^{(m)}-Y_t^{(n)}|^2dt\rightarrow0\condition{$(n\rightarrow\infty)$}
\end{equation*}
Ta chứng minh
\begin{equation*}
	A\equiv\int_0^t\sigma(s,Y_s^{(n)}dB_s\xrightarrow{L^2(P)}\int_0^t\sigma(s,X_s)dB_s:=B
\end{equation*}
Thật vậy:
\begin{equation*}
	\begin{split}
		||A-B||^2_{L^2(P)}&=E\left[(\int_0^t(\sigma(s,Y_s^{(n)})-\sigma(s,X_s))dB_s)^2\right]\\
		&=E\int_0^t|\sigma(s,Y_s^{(n)})-\sigma(s,X_s)|^2ds
	\end{split}
\end{equation*}
theo phép đẳng cự Itô:
\begin{equation*}
	E\left|\int_0^tf(s,\omega)dB_s \right|^2=E\int_0^t|f(s,\omega)|^2ds
\end{equation*}
mà ta lại có $Y_s^{(n)}$ hội tụ đều với $X_s$ hầu chắc chắn, do đó vế trái $\rightarrow0$. Ta cũng có:
\begin{equation*}
	\int_0^tb(s,Y_s^{(n)})ds\xrightarrow{L^2(P)}\int_0^tb(s,X_s)ds
\end{equation*}
vì theo bất đẳng thức Hölder $(\int fg)^2\leq\int f^2 \int g^2$ áp dụng cho $g=1:(\int_0^t f)^2\leq t\int_0^t f^2$, ở đây ta có:
\begin{equation*}
	\left|\int_0^t(b(s,Y_s^{(n)})-b(s,X_s))ds \right|^2\leq t.\int_0^t|b(s,Y_s^{(n)})-b(s,X_s)|^2ds
\end{equation*}
cho nên
\begin{equation*}
	E\left|\int_0^t(b(s,Y_s^{(n)}-b(s,X_s) \right|^2\leq tE.\int_0^t|\underbrace{ b(s,Y_s^{(n)})-b(s,X_s)}_{\text{dần đến 0}}|^2ds
\end{equation*}
Do đó, qua giới hạn trong $L^2(P)$ cả hai vế của \eqref{eq:stt5}, ta nhận thấy $X_t$ thoả mãn phương trình (PT) đã cho:
\begin{equation*}
	\begin{cases}
		dX_t=b(t,X_t)dt+\sigma(t,X_t)dB_t\\
		X_0=Z
	\end{cases}\condition{(PT)}
\end{equation*}
Điều đó hoàn tất chứng minh Định lý tồn tại và duy nhất lời giải của phương trình vi phân ngẫu nhiên.
\subsection{Hướng đi khác cho điều kiện tồn tại và duy nhất nghiệm}
Ta xét phương trình vi phân ngẫu nhiên trong các điều kiện sau:
\begin{enumerate}[(1)]
	\item Không gian xác suất $(\Omega,\mathscr{F},P)$ với $\sigma$ - trường $\{\mathscr{F}_t,i\in [0,T]\}$
	\item Quá trình Wiener $m$ - chiều $W(t)=\left\{W^{(1)}(t),W^{(2)},...,W^{(m)}(t) \right\}$ thích nghi với họ $\sigma$ - trường $\{\mathscr{F}_t,t\in [0,T] \}$ (điều đó có nghĩa là $W(0)=0;\forall t\in[0,T]$ đại lượng $W(t)\mathscr{F}_t$ - đo được và số gia $W(t+s)-W(t)$ khi $s\geq0$ không phụ thuộc vào $\sigma$ - trường $\mathscr{F}_t$).
	\item $\xi_0$ là vectơ ngẫu nhiên $\mathscr{F}_0$ - đo được và vectơ $\xi_0$ không phụ thuộc vào quá trình $W(t)$.
	\item Các hàm số $a(t,x)$ và $\delta(t,x)\left(t\in [0,T];x\in R^m \right)$ nhận các giá trị trong $R^m$ và $\mathscr{L}(R^m)$ tương ứng, với $\mathscr{L}(R^m)$ là tập các toán tự tuyến tính từ $R^m$.\\
\end{enumerate}
Cụ thể với phương trình như sau:
\begin{equation}\label{eq:c3e1}
	d\xi(t)=a(t,\xi(t))dt+\delta(t,\xi(t))dW(t)
\end{equation}
với điều kiện ban đầu: $\xi(0)=\xi_0$\\
Phương trình \eqref{eq:c3e1} có thể viết dưới dạng tích phân:
\begin{equation*}
	\xi(t)=\xi_0+\int_{0}^{t}a(s,\xi(s))ds+\int_{0}^{t}\delta(s,\xi(s))dW(s) \condition{$t\in[0,T]$}
\end{equation*}
trong đó $\xi(t)$ là quá trình ngẫu nhiên cần tìm. Nghiệm của phương trình \eqref{eq:c3e1} với điều kiện ban đầu $\xi(0)=\xi_0$ là quá trình ngẫu nhiên m - chiều $\xi(t),t\in[0,T]$ sao cho:\\
\begin{enumerate}[a)]
	\item Quá trình $\xi(t)$ đo được dần đối với họ $\sigma$ - trường $\{\mathscr{F}_t,t\in[0,T]\}$
	\item Mọi toạ độ của vectơ ngẫu nhiên $\left\{a(t,\xi(t)),t\in[0,T] \right\}$ khả tích tuyệt đối trên đoạn $[0,T]$ với xác suất 1.
	\item Mọi phần tử của ma trận ngẫu nhiên $\left\{\delta(t,\xi(t)),t\in[0,T] \right\}$ bình phương khả tích trên đoạn $[0,T]$ với xác suất 1.
	\item Quá trình $\xi(t)$ có vi phân ngẫu nhiên xác định bởi 
\begin{center}
	$d\xi(t)=a(t,\xi(t))dt+\delta(t,\xi(t))dW(t)$ và $\xi(0)=\xi_0$
\end{center}
\end{enumerate}
Chú ý rằng khi $\delta(t,x)=0$ thì phương trình \eqref{eq:c3e1} sẽ là phương trình vi phân thường, thoả điều kiện ngẫu nhiên ban đầu. Như vậy ta sẽ giải phương trình vi phân thường với mỗi $\omega$ xác định.\\
Ta nói rằng phương trình \eqref{eq:c3e1} có nghiệm duy nhất nếu đối với 2 nghiệm bất kỳ $\xi_1(t),\xi_2(t)$ ta có hệ thức:\\
\begin{equation*}
	P\left\{\sup_{0\leq t\leq T}|\xi_1(t)-\xi_2(t)|>0 \right\}=0
\end{equation*}
\subsubsection{Định lý tồn tại và duy nhất nghiệm}
Giả sử các hệ số của phương trình \eqref{eq:c3e1} thoả các điều kiện:
\begin{enumerate}[(1)]
	\item Với mọi $t\in[0,T],x\in R^m$ ta có:
\begin{equation*}
	|a(t,x)|^2+|\delta(t,x)|^2\leq k(1+|x|^2)
\end{equation*}
trong đó k là một hằng số; còn $|\delta(t,x)|^2=\sum\limits_{i,j=1}^m(\delta^{ij}(t,x))^2\delta^{ij}(t,x)$ là các phần tử của ma trận $\delta(t,x)$
	\item Với mọi $R>0$, tồn tại hằng số $c_R$ sao cho, với $|x|\leq R;y\leq R$ và $t\in[0,T]$ ta sẽ có:
\begin{equation*}
	|a(t,x)-a(t,y)|^2+|\delta(t,x)-\delta(t,y)|^2\leq c_R|x-y|^2
\end{equation*}
Khi đó sẽ tồn tại duy nhất nghiệm liên tục $\xi(t),t\in[0,T]$ của phương trình \eqref{eq:c3e1}. Chứng minh định lý trong quyển \textbf{Stochastic Differential Equations - An Introduction with Application} của B. Oksendal.
\end{enumerate}
Mặc khác ta thấy rằng nghiệm của phương trình \eqref{eq:c3e1} có thể xác định bằng phương pháp xấp xỉ liên tiếp như sau:\\
Giả sử các hệ số $a(t,x),\delta(t,x)$ thoả điều kiện (1) của định lý 3.1.1 và điều kiện (2*) tồn tại hằng số $c>0$ sao cho $\forall t\in[0,T],x,y\in R^m$
\begin{equation*}
	|a(t,x)-a(t,y)|^2+|\delta(t,x)-\delta(t,y)|^2\leq c|x-y|^2
\end{equation*}
Ngoài ra ta còn giả định rằng $M|\xi_0|^2<\infty$. Khi đó, đặt $\eta_0(t)=\xi_0$ và:
\begin{equation*}
	\eta_{n+1}(t)=\xi_0+\int_{0}^ta(s,\eta_n(s))ds+\int_0^t\delta(s,\eta_n(s))dW(s) \condition{n=1,2,...}
\end{equation*}
Sau khi bộ sung thêm những mệnh đề phụ trợ ta sẽ có được bất đẳng thức:
\begin{equation*}
	E\sup_{0\leq t\leq T}|\eta_{n+1}(t)-\eta_n(t)|^2\leq\frac{M^nT^n}{n!} \condition{n=0,1,2,...}
\end{equation*}
trong đó M là một hằng số nào đó.\\3.1
Các bất đẳng thức trên cho ta kết luận được rằng chuỗi:
\begin{equation*}
	\eta_0(t)+\sum_{n=0}^\infty(\eta_{n+1}(t)-\eta_n(t))
\end{equation*}
hội tụ đều với xác suất 1, nó chính là nghiệm của phương trình \eqref{eq:c3e1} khi các hệ số của phương trình thoả các điều kiện (1), (2) và điều kiện ban đầu $M|\xi_0|^2<\infty$.\\
Nếu các điều kiện của định lý 3.1 được thoả và $M|\xi_0|^{2p}<\infty$ với một số p nguyên bất kỳ thì nghiệm của phương trình \eqref{eq:c3e1} với điều kiện ban đầu $\xi_p$ thoả điều kiện $E|\xi(t)|^{2p}\leq k_p(1+E|\xi_0|^{2p})$ với $t\in[0,T]$ và $E|\xi(t)-\xi_0|^{2p}\leq k'_p(1+E|\xi_0|^{2p})$ với $t\in[0,T]$, trong đó $k_p,k'_p$ là những hằng số chỉ phụ thuộc vào $p,k$ và $T$.\\
Trong các điều kiện của định lý 3.1, nghiệm $\xi(t),t\in[0,T]$ của phương trình \eqref{eq:c3e1} có tính chất Markov đối với họ $\sigma$ - trường ${\mathscr{F}_t,t\in[0,T]}$. Điều đó đồng nghĩa với mọi $0\leq s\leq t\leq R,G\in \mathscr{B},\mathscr{B}$ là $\sigma$ - trường các tập Borel của $R^m$, với xác suất 1 ta sẽ có hệ thức:
\begin{equation*}
	P\{\xi(t)\in G/\mathscr{F}_s\}=P\{\xi(t)\in G/\xi(\mathscr{F})\}
\end{equation*}
Như vậy quá trình ngẫu nhiên $\xi(t),t\in[0,T]$ là một hàm ngẫu nhiên Markov với phân phối ban đầu:
\begin{equation*}
	\mu(G)=P\{\xi_0\in G\},G\in \mathscr{B}
\end{equation*}
Xác suất chuyển $p(s,x,t,G)$ của hàm ngẫu nhiên Markov $\xi(t)$ xác định bởi công thức:
\begin{equation*}
	p(s,x,t,G)=P\{\xi_{sx}(t)\in G\},0\leq s\leq t\leq T \condition{$x\in R^m,G\in\mathscr{B}$}
\end{equation*}
trong đó $\xi_{sx}(t)$ là nghiệm của phương trình:
\begin{equation}\label{eq:c3e2}
	\xi_{sx}(t)=x+\int_{s}^{t}a(\tau,\xi_{sx}(\tau))d\tau+\int_{s}^{t}\delta(\tau,\xi_{sx}(\tau))dW(\tau)
\end{equation}
Trong \eqref{eq:c3e2}, x là vectơ ngẫu nhiên trong $R^m,t\in[s,T]$\\
\subsubsection{Định lý về quá trình khuếch tán và ma trận khuếch tán}
Giả sử các hàm số $a(t,x)$ và $\delta(t,x)$ thoả các điều kiện của định lý 3.1, ngoài ra chúng là những hàm liên tục theo tập hợp các biến số. Khi đó quá trình ngẫu nhiên $\xi(t),t\in[0,T]$ là nghiệm của phương trình \eqref{eq:c3e1} sẽ là một quá trình khuếch tán với vectơ chuyển dịch $a(t,x)$ và ma trận khuếch tán:
\begin{equation*}
	B(t,x)=\delta(t,x)\delta^T(t,x)
\end{equation*}
Như vậy, lý thuyết về quá trình vi phân ngẫu nhiên cho ta khả năng xây dựng các quá trình khuếch tán trong những điều kiện khá rộng đối với các hệ số $a(t,x)$ và $B(t,x)$.\\
Ngoài ra, nếu bổ sung thêm về tính trơn của các hàm $a(t,x),\delta(t,x)$ ta có thể chứng minh sự tồn tại đạo hàm liên tục của hàm số:
\begin{equation*}
	u(s,x)=E\mathscr{f}(\xi_{sx}(t)),0\leq s\leq t\leq T,x\in R^m
\end{equation*}
theo biến x, khi đó ta sẽ thu được phương trình ngược Kolmogorov.\\
\subsubsection{Định lý về hàm liên tục và khả vi hai lần}
Giả sử các hàm số $a(t,x)$ và $\delta(t,x)$ thoả điều kiện của định lý 3.1, chúng là những hàm liên tục và khả vi 2 lần theo biến $x$. Ngoài ra, với những số xác định $p>0$ và $k>0$ ta có bất đẳng thức:
\begin{dmath*}
	\sum\limits_{t,k=1}^{m}\left|\dfrac{\partial a^i(t,x)}{\partial x^k}\right|+\sum\limits_{i,j,k=1}^{m}\left|\dfrac{\partial^2a^i(t,x)}{\partial x^i\partial x^k}\right|+\sum_{i,j,k=1}^{m}\left|\dfrac{\partial\delta^{ij}(t,x)}{\partial x^k} \right|+\sum_{i,j,k,l=1}^{m}\left|\dfrac{\partial^2\delta^{ij}(t,x)}{\partial x^k\partial x^l} \right|\leq k(1+|x|^p)	
\end{dmath*}
khi đó nếu hàm $f(x),(x\in R^m)$ khả vi liên tục hai lần với các giá trị thực sao cho:
\begin{equation*}
	\left|f(x)\right|+\sum_{i=1}^{m}\left|\dfrac{\partial f(x)}{\partial x^2} \right|+\sum_{i,k=1}^{m}\left|\dfrac{\partial^2f(x)}{\partial x^i \partial x^k}\leq k(1+|x|^p \right| \condition{p>0}
\end{equation*}
ta sẽ có hàm:
\begin{equation*}
	u(s,x)=Ef(\xi_{sx}(t)),0\leq s\leq t \leq T,x\in R^m
\end{equation*}
trong đó $\xi_{sx}(t)$ là nghiệm của phương trình \eqref{eq:c3e2}, hai lần khả vi theo x, khả vi liên tục theo s và thoả phương trình:
\begin{equation*}
	\dfrac{\partial u(s,x)}{\partial s}+\sum_{i=1}^{m}a^i(s,x)\dfrac{\partial u(s,x)}{\partial x^2}+\dfrac{1}{2}\sum_{i,j,k=1}^{m}\delta^{ij}(s,x)\delta^{kj}(s,x)\dfrac{\partial^2u(s,x)}{\partial x^i\partial x^k}=0
\end{equation*}
trong miền $s\in(0,t),x\in R^m$ với điều kiện ban đầu là $\lim\limits_{s\rightarrow t}u(s,x)=f(x)$.\\
Hệ quả của định lý trên là sự tồn tại và duy nhất nghiệm của bài toán Cauchy đối với phương trình đạo hàm riêng dạng parabolic.\\
\subsubsection{Định lý về nghiệm duy nhất với mọi hàm khả vi liên tục hai lần}
Giả sử trong miền $0\leq s<T;x\in R^m$, xác định một toán tử vi phân
\begin{equation*}
	\mathscr{L}u(s,x)=\dfrac{\partial u}{\partial s}(s,x)+\sum_{i=1}^{m}a^i(s,x)\dfrac{\partial u(s,x)}{\partial x^i}+\dfrac{1}{2}\sum_{i,j=1}^{m}b^{ij}(s,x)\dfrac{\partial^2u(s,x)}{\partial x^i\partial x^j}
\end{equation*}
dạng parabolic. Nếu ma trận $B(t,x)$ với các phần tử $b^{ij}(t,x);i,j=1,2,...,m$ sao cho $B(t,x)=\delta(t,x)\delta^T(t,x)$ và các hàm $a(t,x),\delta(t,x)$ thoả điều kiện của định lý 3.3 thì khi đó bài toán Cauchy với:
\begin{equation*}
	\mathscr{L}u(s,x)=0,\lim\limits_{s\rightarrow T}u(s,x)=f(x)
\end{equation*} 
sẽ có nghiệm duy nhất với mọi hàm khả vi liên tục hai lần $f(x)$ đồng thời nó cùng với các đạo hàm riêng cho đến bậc hai của nó tăng đến vô cực không nhanh hơn một bậc nào bất kỳ của $|x|$. Trong trường hợp này, nghiệm $u(s,x)$ của bài toán Cauchy sẽ được biểu diễn dưới dạng:
\begin{equation*}
	u(s,x)=Ef(\xi_{sx}(t)),0\leq s<T,x\in R^m
\end{equation*}
với $\xi_{sx}(t),(t\in[s,T])$ là nghiệm của phương trình \eqref{eq:c3e2}\\
Định lý này chỉ ra rằng khi nghiên cứu về phương trình đạo hàm riêng dạng parabolic ta có thể sử dụng các kết quả của lý thuyết về phương trình vi phân ngẫu nhiên. Đặc biệt trong định lý 3.4 không đòi hỏi tính không suy biến của ma trận $B(t,x)$.\\
\subsubsection{Định nghĩa nghiệm mạnh và nghiệm yếu}
Nghiệm đã xác định trong định lý 3.1 được gọi là một nghiệm mạnh vì các cơ sở sau đã được xác định trước
\begin{enumerate}[1)]
	\item Không gian xác suất cơ bản $(\Omega,\mathscr{F},P)$
	\item Họ các $\sigma$ - trường con đầy đủ của $\mathscr{F}:\{\mathscr{F}_t,t\in[0,1]\}$
	\item Quá trình Wiener $\{W_t,t\in[0,T]\}$ đã xác định sao cho $\{W_t,\mathscr{F}_t,t\in[0,T]\}$ lập thành Martingan.
\end{enumerate}
Khi chỉ cho trước các hàm $a(t,x)$ và $\delta(t,x)$ mà không cho trước ba yếu tố trên ta sẽ có nghiệm theo nghĩa yếu. Hay nói cách khác, phương trình:
\begin{equation*}
	d\xi_t=a(t,\xi)dt+\delta(t,\xi)dW(t)
\end{equation*}
được gọi là có nghiệm yếu nếu tìm được
\begin{enumerate}[(1*)]
	\item Không gian xác suất $(\Omega^*,\mathscr{F}^*,P^*)$
	\item Họ các $\sigma$ - trường con đầy đủ của $\mathscr{F}^*:\{\mathscr{F}^*_t,t\in[0,1]\}$
	\item Quá trình Wiener $\{W_t^*,t\in[0,T]\}$ sao cho $\{W_t^*,\mathscr{F}^8_t,t\in[0,T]\}$ lập thành Martingan.
	\item $\{\xi_t^*,\mathscr{F}_t^*,t\in[0,T]\}$ là một quá trình liên tục và thích nghi (tương thích) sao cho:
	\begin{equation*}
		d\xi_t^*=a(t,\xi)dt+\delta(t,\xi)dW(t)	
	\end{equation*}
\end{enumerate}
Như vậy ta nói rằng $\{\xi_t^*,t\in[0,T]\}$ là nghiệm yếu của phương trình vi phân ngẫu nhiên đã cho. Điều kiện ban đầu của nghiệm yếu là hàm phân phối xác suất $F$ cho trước, ta phải tìm nghiệm yếu sao cho $\xi_0^*$ có hàm phân phối bằng $F$.\\
\textbf{Chú ý:}
\begin{itemize}[*]
	\item Tính duy nhất của nghiệm mạnh được hiểu theo nghĩa có cùng quỹ đạo.
	\item Tính duy nhất của nghiệm yếu được hiểu theo nghĩa có cùng phân phối xác suất.
	\item Tên gọi nghiệm mạnh và nghiệm yếu được đặt ra vì nghiệm mạnh sẽ là nghiệm yếu, nhưng nghiệm yếu không thể trở thành nghiệm mạnh.
\end{itemize}
\pagebreak
\section{Dạng của phương trình vi phân ngẫu nhiên:}

\subsection{Một số phương trình đặc biệt}
\subsubsection{Phương trình vi phân tuyến tính thuần nhất}
Phương trình vi phân ngẫu nhiên tuyến tính thuần nhất là phương trình có dạng:
\begin{equation}\label{eq:c3.2.1}
	d\xi(t)=\alpha(t)\xi(t)dt+\beta(t)\xi(t)dW_t
\end{equation} 
trong đó $\alpha(t),\beta(t)$ là các hàm số của t và $W_t$ là quá trình Wiener với điều kiện ban đầu $\xi(0)=\xi_0$\\
\textbf{Cách giải:} Trước tiên ta tìm phương pháp giải \eqref{eq:c3.2.1} trong trường hợp đặc biệt khi $\alpha(t)=0$. Cụ thể xét các ví dụ sau:\\
a) Giải phương trình: \inlineequation[eq:c3.2.2]{\begin{cases} d\xi(t)=\beta(t)\xi(t)dW_t \\ \xi(0)=1 \end{cases}}\\
Ta xét quá trình: \inlineequation[eq:c3.2.3]{\eta(t):=-\dfrac{1}{2}\int_{0}^{t}\beta^2(s)ds+\int_{0}^{t}\beta(s)dW_s} với vi phân tương ứng là: $d\eta=-\dfrac{1}{2}\beta^2(t)dt+\beta(t)dW_t$.\\
Sử dụng công thức Itô với $\varphi(t,x)=e^x$ ta sẽ có:
\begin{dmath*}
d\xi(t)=d\varphi(t,\eta_t)=\dfrac{\partial\varphi}{\partial x}d\eta+\dfrac{1}{2}\dfrac{\partial\varphi^2}{\partial x^2}\beta^2(t)dt=e^\eta(-\dfrac{1}{2}\beta^2(t)dt+\beta(t)dW_t+\dfrac{1}{2}\beta^2(t)dt)=\beta(t)\xi(t)dW_t
\end{dmath*}
Từ đó, theo \eqref{eq:c3.2.3} ta sẽ có nghiệm của phương trình \eqref{eq:c3.2.2} là:
\begin{equation*}
	\xi(t)=\exp\left\{-\dfrac{1}{2}\int_{0}^{t}\beta^2(s)ds+\int_{0}^{t}\beta(s)sW_s \right\}
\end{equation*}
Nếu ta thay điều kiện ban đầu $\xi(0)=1$ bởi điều kiện $\xi(0)=\xi_0$ ta sẽ có nghiệm của phương trình: \inlineequation[eq:c3.2.4]{\begin{cases}
	d\xi(t)=\beta(t)\xi(t)dW_t\\
	\xi(0)=\xi_0
\end{cases}} là quá trình ngẫu nhiên:
\begin{equation}\label{eq:c3.2.5}
	\xi(t)=\xi_0\exp\left\{\int_{0}^{t}\beta(s)dW_s-\dfrac{1}{2}\int_{0}^{t}\beta^2(s)ds \right\}
\end{equation}
b) Giải phương trình
\begin{equation}\label{eq:c3.2.6}
\begin{cases}
	d\xi(t)=\alpha(t)\xi(t)dt+\beta(t)\xi(t)dt\\
	\xi(0)=\xi_0
\end{cases}
\end{equation}
Ta tìm nghiệm của \eqref{eq:c3.2.6} dưới dạng: $\xi(t)=\xi_1(t).\xi_2(t)$ trong đó $\xi_1(t)$ thoả điều kiện: \inlineequation[eq:c3.2.7]{\begin{cases}
	d\xi_1(t)=\beta(t)\xi_1(t)dW_t\\
	\xi_1(0)=1
\end{cases}} và $\xi_2(t)$ thoả điều kiện: 
\begin{equation}\label{eq:c3.2.8}
\begin{cases}
	d\xi_2(t)=A(t)dt+B(t)dW_t\\
	\xi_2(0)=1
\end{cases}
\end{equation} với $A(t)$ và $B(t)$ là những hàm ta sẽ chọn sau này.\\
Khi đó ta có:
\begin{dmath*}
d\xi(t)=d(\xi_1,\xi_2)=\xi_1d\xi_2+\xi_2d\xi_1+\beta(t)\xi_1B(t)dt=\beta(t)\xi(t)dW_t+\xi_1d\xi_2+\beta(t)\xi_1B(t)dt	
\end{dmath*}
Ta chọn $A(t)$ và $B(t)$ sao cho:
\begin{equation*}
	d\xi_2(t)+\beta(t)B(t)dt=\alpha(t)\xi_2(t)dt
\end{equation*}
Cụ thể ta lấy: $A(t)=\alpha(t)\xi_2(t)$ và $B(t)=0$.\\
Từ hệ thức \eqref{eq:c3.2.8} sẽ cho ta:
\begin{equation}\label{eq:c3.2.9}
	\begin{cases}
		d\xi_2(t)=\alpha(t)\xi_2(t)dt\\
		\xi_2(0)=1
	\end{cases}
\end{equation}
Từ \eqref{eq:c3.2.9} sẽ cho ta:
\begin{equation*}
	\xi_2(t)=\exp\left(\int_{0}^{t}\alpha(s)ds \right)
\end{equation*}
Mặt khác theo phần a), phương trình \eqref{eq:c3.2.7} cho ta nghiệm là:
\begin{equation*}
	\xi_1(t)=\xi_0\exp\left\{\int_{0}^{t}\beta(s)dW_s-\dfrac{1}{2}\int_{0}^{t}\beta^2(s)ds \right\}
\end{equation*}
Kết hợp $\xi_1(t)$ và $\xi_2(t)$ ta có nghiệm của phương trình \eqref{eq:c3.2.6} là:
\begin{equation}\label{eq:c3.2.10}
	\xi(t)=\xi_1(t)\xi_2(t)=\exp\left\{\int_{0}^{t}\beta(s)dW_s+\int_{0}^{t}\left[\alpha(s)-\dfrac{1}{2}\beta^2(s) \right]ds\right\}
\end{equation}
\subsubsection{Phương trình Black-Scholes}
Phương trình mô tả sự biến động của giá cổ phiếu $S_t$ theo thời gian t:
\begin{equation}\label{eq:c3.2.11}
	dS_t=\alpha S_tdt+\beta S_tdW_t
\end{equation}
trong đó $\alpha$ và $\beta$ là những hằng số.\\
Đây là trường hợp đặt biệt của phương trình dạng \eqref{eq:c3.2.1}. Do đó từ \eqref{eq:c3.2.10} ta suy ra nghiệm của nó sẽ là:
\begin{equation}\label{eq:c3.2.12}
	S_t=\exp\left[\left(\alpha-\dfrac{\beta^2}{2} \right)t+\beta W_t \right]
\end{equation}
\subsubsection{Phương trình mũ Itô}
Dạng của phương trình như sau:
\begin{equation}\label{eq:c3.2.13}
	\begin{cases}
		dX(t)=\dfrac{1}{2}X(t)dt+X(t)dW_t\\
		X(0)=X_0
	\end{cases}
\end{equation}
Từ \eqref{eq:c3.2.10} ta suy ra nghiệm của nó sẽ là: $X(t)=X_0exp(W_t)$\\
\subsection{Phương trình vi phân tuyến tính dạng tổng quát}
Phương trình vi phân tuyến tính là phương trình có dạng
\begin{equation}\label{eq:c3.2.14}
	d\xi(t)=(\alpha(t)\xi(t)+f(t))dt+(\beta(t)\xi(t)+g(t))dW_t
\end{equation}
trong đó $\alpha(t),\beta(t),f(t),g(t)$ là những hàm của t và $W_t$ là quá trình Wiener với điều kiện ban đầu $\xi(0)=\xi_0$.\\
\textbf{Cách giải:} Ta tìm nghiệm của phương trình \eqref{eq:c3.2.14} dưới dạng: $\xi(t)=\xi_1(t)\xi_2(t)$, trong đó $\xi_1(t)$ là nghiệm của phương trình tuyến tính thuần nhất tương ứng.
\begin{equation}\label{eq:c3.2.15}
	\begin{cases}
		d\xi_1(t)=\alpha(t)\xi_1(t)dt+\beta(t)\xi_1(t)dW_t\\
		\xi_1(0)=1
	\end{cases}
\end{equation}
và $\xi_2(t)$ là nghiệm của phương trình:
\begin{equation*}
	\begin{cases}
		d\xi_2(t)=A(t)dt+B(t)dW_t\\
		\xi_2(0)=\xi_0
	\end{cases}
\end{equation*}
Hàm $A(t)$ và $B(t)$ ta sẽ chọn sau này. Khi đó ta có:
\begin{equation*}
	d\xi(t)=d(\xi_1(t)\xi_2(t))=\xi_1d\xi_2+\xi_2d\xi_1+\beta(t)\xi_1B(t)dt
\end{equation*}
Cụ thể ta chọn $A(t)$ và $B(t)$ sao cho:
\begin{equation*}
	\xi_1(t)[A(t)dt+B(t)dW_t]+\beta(t)\xi_1(t)B(t)dt=f(t)dt+g(t)dW_t
\end{equation*}
ta sẽ thu được: \inlineequation[eq:c3.2.16]{\begin{cases}
	A(t):=[f(t)-\beta(t)g(t)]\dfrac
	{1}{\xi_1(t)}\\
	B(t):=g(t)\dfrac{1}{\xi_1(t)}
\end{cases}}\\
Sử dụng kết quả ở phần b), theo công thức \eqref{eq:c3.2.14} ta có:
\begin{equation*}
	\xi_1(t)=\exp\left\{\int_{0}^{t}\beta(s)dW_s+\int_{0}^{t}\left[\alpha(s)-\dfrac{1}{2}\beta^2(s) \right]ds \right\}
\end{equation*}
Theo \eqref{eq:c3.2.16} ta có:
\begin{equation*}
	\xi_2(t)=\xi_0+\int_{0}^{t}[f(s)-\beta(s)g(s)](\xi_1(s))^{-1}ds+\int_{0}^{t}g(s)(\xi_1(s))^{-1}dW_s
\end{equation*}
Kết hợp $\xi_1(t)$ và $\xi_2(t)$ ta có nghiệm của phương trình \eqref{eq:c3.2.14} sẽ là:
\begin{equation*}
	\xi(t)=\xi_1(t)\xi_2(t)
\end{equation*}
\subsection{Phương trình Langevin}
Phương trình Langevin là phương trình có dạng:
\begin{equation}\label{eq:c3.2.17}
	\begin{cases}
		d\xi(t)=-b\xi(t)dt+\delta dW_t\\
		\xi(0)=\xi_0
	\end{cases}
\end{equation}
trong đó $b,\delta$ là những hằng số thực.\\
Phương trình Langevin là một phương trình tuyến tính, do đó áp dụng kết quả phần II, theo công thức \eqref{eq:c3.2.17} ta thu được nghiệm của nó là phương trình ngẫu nhiên:
\begin{equation}\label{eq:c3.2.18}
	\xi(t)=e^{-bt}\left[\xi_0+\int{0}^{t}\delta e^{bs}dW_s \right] \condition{$t\geq 0$}
\end{equation}
Khi đó quá trình xác định bởi \eqref{eq:c3.2.18} được gọi là quá trình Ornstein-Unlenbeck.\\
\subsection{Phương trình Langevin mở rộng}
Phương trình Langevin mở rộng là phương trình có dạng:
\begin{equation}\label{eq:c3.2.19}
	\begin{cases}
		d\xi(t)=\alpha(t)\xi(t)dt+\delta(t)dW_t\\
		\xi(0)=\xi_0
	\end{cases}
\end{equation}
trong đó $\alpha(t)$ là một quá trình liên tục thích nghi.\\
Từ lời giải của phương trình vi phân tuyến tính tổng quát ta sẽ có nghiệm của phương trình Langevin mở rộng là:
\begin{equation}\label{eq:c3.2.20}
	\xi(t)=\exp\left(\int_{0}^{t}\alpha(s)ds\right)\left[\xi_0+\int_{0}^{t}\exp\left(-\int_{0}^{r}\alpha(s)ds \right)\delta(r)dW_r \right]
\end{equation}
Xét phương trình tuyến tính có dạng:
\begin{equation}\label{eq:c3.2.21}
	\begin{cases}
		d\xi(t)=\alpha(t)\xi(t)+f(t)]dt+\sum\limits_{i=1}^{m}\left[\beta^i(t)\xi(t)+g^i(t)\right]dW_t^i\\
		\xi(0)=\xi_0
	\end{cases}
\end{equation}
Khi đó nghiệm của nó sẽ là:
\begin{dmath}\label{eq:c3.2.22}
\xi(t)=G(t)\left[\xi_0+\int_{0}^{t}\left(G(s)\right)^{-1}\left(\alpha(s)-\sum\beta^i(s)g^i(s)\right)ds \right]+\int_{0}^{t}\sum\limits_{i=1}^{m}(G(s))^{-1}g^i(s)dW_t^i	
\end{dmath}
trong đó:
\begin{equation*}
	G(t):=\exp\left(\int_{0}^{t}\left[\alpha(s)-\sum_{i=1}^{m}\dfrac{(\beta^i)^2}{2} \right]ds+\int_{0}^{t}\sum_{i=1}^{m}\beta^i(s)dW_s^i \right)
\end{equation*}
\pagebreak
\section{Phương trình vi phân tuyến tính Itô - Levy}
\subsection{Định nghĩa phương trình vi phân tuyến tính Itô - Levy}
\subsubsection{Định nghĩa}
Phương trình vi phân tuyến tính Itô - Levy là phương trình có dạng:
\begin{dmath}\label{eq:dn1}
	dX(t)=[\alpha(t,\omega)X(t^-)+A(t,\omega)]dt+[\beta(t,\omega)X(t^-)+B(t,\omega)]dW(t)\\
	+\int_{(R_0)^{n_2}}[\gamma(t,z,\omega)X(t^-)+G(t,z,\omega)]\overline{N}(dt,dz)
\end{dmath}
Trong đó điều kiện ban đầu $X(0)=x_0$, với:
\begin{center}
$\alpha(t,\omega); \beta(t,\omega);A(t,\omega);B(t,\omega);\gamma(t,z,\omega);G(t,z,\omega);$\\
$\forall t \geqslant 0; z \in R_0; \omega \in \Omega$
\end{center}
là những quá trình ngẫu nhiên khả đoán cho trước với:
\begin{center}
	$\gamma(t,z,\omega)>-1;\forall (t,z,\omega)\in [0,\inf\}\times R_0 \times \Omega$
\end{center}
và thoả các điều kiện:
\begin{center}
	$\int\limits_{0}^t[|\alpha(s,\omega)|-\dfrac{1}{2}\beta^2(s,\omega)+\int\limits_{R_0} \gamma^2(s,z,\omega)v(dz)]ds<\infty$; h.c\\
	$\int\limits_{0}^t [|A(s,\omega)-\dfrac{1}{2}B^2(s,\omega)+\int\limits_{R_0}G^2(s,z,\omega)v(dz)]ds<\infty$; h.c
\end{center}
khi $A(t,\omega)\equiv B(t,\omega)\equiv G(t,z,\omega)\equiv 0$, h.c; ta gọi đó là quá trình vi phân ngẫu nhiên tuyến tính thuần nhất hay còn gọi là phương trình vi phân ngẫu nhiên hình học.\\
\subsubsection{Định lý giải phương trình vi phân ngẫu nhiên}
Cho phương trình vi phân ngẫu nhiên tuyến tính Itô - Levy nêu trong định nghĩa trên, khi đó phương trình \eqref{eq:dn1} sẽ có nghiệm là:
\begin{dmath}\label{eq:dl1}
	\dfrac{X(t)}{X_1(t^-)}=\left\{x_0+\int_{R_0}\dfrac{1}{X_1(s^-)}[A(s,\omega)-\beta(s,\omega)B(s,\omega)-\int_{R_0} \dfrac{\gamma(s,z,\omega)G(s,z,\omega)}{1+\gamma(s,z,\omega)}v(dz)]ds+\int_{0}^{t} \dfrac{B(s,\omega)}{X_1(s)}+\int_{0}^{t} \int_{R_0} \dfrac{G(s,z,\omega)}{X_1(s^-)(1+\gamma(s,z,\omega))} \overline{N}(ds,dz) \right\}
\end{dmath}
Trong đó:
\begin{dmath}\label{eq:pt1}
	X_1(t)=exp\left\{\int_{0}^{t}\left[\alpha(s,\omega)-\dfrac{1}{2}\beta^2(s,\omega)+\int_{R_0}\log(1+\gamma(s,z,\omega))-\gamma(s,z,\omega)v(dz) \right]ds+\int_{0}^{1}\beta(s,\omega)dW(s)+\int_{0}^{1}\int_{R_0}\log(1+\gamma(s,t,\omega))\overline{N}(ds,dz) \right\}
\end{dmath}
Để chứng minh định lý trên trước hết ta cần chứng minh bổ đề sau:\\
\subsubsection{Bổ đề}
Cho phương trình vi phân ngẫu nhiên tuyến tính thuần nhất, nghĩa là phương trình có dạng:
\begin{dmath}\label{eq:pt2}
	\dfrac{dX_1(t)}{X_1(t^-)}=\left[\alpha(t,\omega)dt+\beta(t,\omega)dW(t)+\int_{R_0}\gamma(t,z,\omega)\overline{N}(dt,dz) \right]	
\end{dmath}
Với điều kiện ban đầu $X(0)=1$, trong đó:
\begin{equation*}
\alpha(t,\omega);\beta(t,\omega);\gamma(t,z,\omega);t\geq0;z\in R_0;\omega\in \Omega	
\end{equation*}
là những quá trình ngẫu nhiên khả đoán cho trước với
\begin{equation*}
	\gamma(t,z,\omega)>-1;\forall (t,z,\omega)\in [0;\infty\}\times R_0 \times \Omega
\end{equation*}
và thoả điều kiện:
\begin{equation*}
	\int\limits_{0}^{t}\left[|\alpha(s,\omega)|-\dfrac{1}{2}\beta^2(s,\omega)+\int_{R_0}\gamma^2(s,z,\omega)v(dz) \right]ds<\infty \condition{h.c}
\end{equation*}
khi đó nghiệm của phương trình \eqref{eq:pt2} sẽ được cho bởi hệ thức \eqref{eq:pt1}.\\
\subsubsection{Chứng minh bổ đề}
Ta xét hàm $X_1(t)=F(t,H(t));t\geq 0$ với $F(t,x)=e^x$ và $H(T)$ xác định bởi:
\begin{dmath*}
H(t)=\int\limits_{0}^{t}\left[\alpha(s,\omega)-\dfrac{1}{2}\beta^2(s,\omega)+\int_{R_0}\log(1+\gamma(s,z,\omega))-\gamma(s,z,\omega)v(dz) \right]ds+\int_{0}^{1}\beta(s,\omega)dW(s)+\int_{0}^{1}\int_{R_0}\log(1+\gamma(s,z,\omega))\overline{N}(ds,dz)	
\end{dmath*}
Áp dụng công thức Itô cho $X_1(t)=F(t,H(t))$, ta sẽ thu được:
\begin{dmath*}
dX_1(t)=e^{H(t)}\left[\left(\alpha(t,\omega)-\dfrac{1}{2}\beta^2(t,\omega)+\int_{R_0}[\log(1+\gamma(t,z,\omega))-\gamma(t,z,\omega)]v(dz) \right)dt\right]+e^{H(t)}\left[\dfrac{1}{2}\beta^2(t,\omega)dt+\beta(t,\omega)dW(t) \right]	+\int_{R_0}e^{H(t)}[\gamma((t,z,\omega)-\log(1+\gamma(t,z,\omega)))]v(dz)dt+\int_{R_0}e^{H(t^-)}\gamma(t,z,\omega)\tilde{N}(dt,dz)=X_1(t^-)\left[\alpha(t,\omega)dt+\beta(t,\omega)dW(t)+\int_{R_0}\gamma(t,z,\omega)\tilde{N}(dt,dz) \right]\blacksquare
\end{dmath*}
\subsubsection{Chứng minh định lý}
Ta tìm nghiệm phương trình \eqref{eq:dn1} bằng phương pháp tách nghiệm, nghĩa là tìm nghiệm của nó dưới dạng tích
\begin{equation}\label{eq:pt3}
	X(t)=X_1(t^-).X_2(t^-)
\end{equation}
trong đó:
\begin{itemize}
	\item $X_1(t)$ là nghiệm của phương trình tuyến tính thuần nhất tương ứng, nghĩa là nó là nghiệm của phương trình \eqref{eq:pt2} nói trong Bổ đề.
	\item $X_2(t)$ là nghiệm của phương trình:
\begin{equation*}
	dX_2(t)=A^*(t,\omega)dt+B^*(t,\omega)dW(t)+\int_{R_0}G^*(t,z,\omega)\tilde{N}(dt,dz)
\end{equation*}
\end{itemize}
Với điều kiện $X_2(0)=x_0$, trong đó $A^*(t,\omega);B^*(t,\omega);G^*(t,z,\omega)$ là những hàm ta sẽ xác định sau.\\
Theo bổ đề, ta có nghiệm $X_1(t^-)$ của phương trình \eqref{eq:pt2} cho bởi hệ thức \eqref{eq:pt1}. Áp dụng hệ quả cho tích $X(t)=X_1(t^-).X_2(t^-)$ nêu trên, ta thu được:
\begin{dmath}
d(X(t))=d(X_1(t^-).X_2(t^-))=X_1(t^-).dX_2(t)+X_2(t^-)dX_1(t)+\beta(t,\omega)X_1(t^-)B^*(t,\omega)dt+\int_{R_0}\gamma(t,z,\omega)X_1(t^-)G^*(t,z,\omega)\tilde{N}(dt,dz)=\alpha(t,\omega)X_1(t^-)X_2(t^-)+\beta(t,\omega)X_1(t^-)X_2(t^-)+\int_{R_0}\gamma(t,z,\omega)X_1(t^-)X_2(t^-)\tilde{N}(dt,dz)+X_1(t^-)A^*(t,\omega)dt+X_1(t^-)B^*(t,\omega)dW(t)+X_1(t^-)\int_{R}G^*(t,z,\omega)\tilde{N}(dt,dz)+\beta(t,\omega)X_1(t^-)B^*(t,\omega)dt+\gamma(t,z,\omega)X_1(t^-)G*(t,z,\omega)\tilde{N}(dt,dz)	
\end{dmath}
Mặt khác, $X(t)$ là nghiệm của phương trình \eqref{eq:dn1}, từ đó so sánh giữa \eqref{eq:dn1} và \eqref{eq:pt3}, ta thu được hệ phương trình:
\begin{dmath*}
\begin{cases}
	A(t,\omega)=X_1(t^-)\left[A^*(t,\omega)+B(t,\omega)B^*(t,\omega)+\int\limits_{R_0}\gamma(t,z,\omega)G(t,z,\omega)v(dz) \right] \\
	B(t,\omega)=X_1(t^-)B^*(t,\omega) \\
	\int\limits_{R_0}G(t,z,\omega)\tilde{N}(dt,dz)=X_1(t^-)\int\limits_{R_0}(1+\gamma(t,z,\omega))G^*(t,z,\omega)\tilde{N}(dt,dz)
\end{cases}	
\end{dmath*}
Suy ra:
\begin{dmath*}
\begin{cases}
	A^*(t,\omega)=\dfrac{1}{X_1(t^-)} \left[A(t,\omega)-B(t,\omega)\beta(t,\omega)-\int_{R_0} \dfrac{\gamma(t,z,\omega)G(t,z,\omega)}{1+\gamma(t,z,\omega)} v(dz) \right] \\
	B^*(t,\omega)=\dfrac{B(t,\omega)}{X_1(t^-)} \\
	G^*(t,z,\omega)=\dfrac{G(t,z,\omega)}{X_1(t^-)(1+\gamma(t,z,\omega)}
\end{cases}	
\end{dmath*}
Đặt $X_1(t)$ cho bởi hệ thức \eqref{eq:pt1}, và các biểu thức của $A^*(t,\omega);B^*(t,\omega);G^*(t,z,\omega)$ đã xác định được vào biểu thức \eqref{eq:pt3}, ta sẽ có nghiệm \eqref{eq:dn1}. $\blacksquare$\\
\section{Một số ứng dụng của phương trình vi phân ngẫu nhiên}
\subsection{Bài toán về chọn danh mục đầu tư rủi ro}
Trước hết, ta mô hình hoá bài toán về tài sản và phương án đầu tư trên thị trường bằng các công cụ của giải tích ngẫu nhiên. Khi xét đến biến động của dòng tài sản trên thị trường, người ta thường xét đến hai loại tài sản: phi rủi ro (như trái phiếu...) và tài sản rủi ro (như cổ phiếu, bất động sản...)\\
Loại tài sản phi rủi ro với biến động giá $X_0(t)$ được xét qua phương trình vi phân:
\begin{equation*}
	dX_0(t)=\lambda(t,\omega)X_0(t)dt \condition{$X_0(0)=1,t\in[0,T]$}
\end{equation*}
Loại tài sản rủi ro với biến động giá $X_1(t)$ thường được xét qua phương trình vi phân ngẫu nhiên Itô có dạng:
\begin{equation}\label{eq:ud1}
	dX_1(t)=\alpha(t,\omega)X_1(t)dt+\beta(t,\omega)X_1(t)dB_t(\omega) \condition{$X_1(0)>0,t\in[0,T]$}
\end{equation}
trong đó: $\lambda(t,\omega)=\lambda(t);\alpha(t,\omega)=\alpha(t);\beta(t,\omega)=\beta(t)$ là những quá trình ngẫu nhiên thoả điều kiện:
\begin{equation*}
	E\int\limits_{0}^{T}[|\lambda(t)|+\alpha(t)+\beta^2(t)]dt<\infty;
\end{equation*}
Ta ký hiệu: $\tau_0(t)=\tau_0(t,\omega);\tau_1(t)=\tau_1(t)(t,\omega);t\in[0,T];\omega\in \Omega$ là những đơn vị vốn đầu tư cho loại tài sản phi rủi ro và rủi ro tương ứng. Khi đó ta gọi: $\tau(t):=(\tau_0(t),\tau_1(t))$ là một phương án đầu tư (một danh mục đầu tư) với tổng giá trị tài sản tại thời điểm t bằng:
\begin{equation}\label{eq:ud2}
	V^\tau(t)=\tau_0(t)X_0(t)+\tau_1(t)X_1(t)	
\end{equation}
Phương án đầu tư được gọi là tự tài trợ (seft-financing) nếu ta có:
\begin{equation}\label{eq:ud3}
	dVT\tau(t)=\tau_0(t)dX_0(t)+\tau_1(t)dX_1(t)
\end{equation}
Từ giả định: $\tau(t);t\in[0,T]$ là tự tài trợ, từ hệ thức \eqref{eq:ud2} ta sẽ có:
\begin{equation*}
	\tau_0(t)=\dfrac{V^\tau(t)-\tau_1(t)X_1(t)}{X_0(t)}
\end{equation*}
Từ các hệ thức trên suy ra:
\begin{equation*}
	dV^\tau(t)=\lambda(t)(V^\tau(t)-\tau_1(t)X_1(t))dt+\tau_1(t)dX_1(t)
\end{equation*}
Sử dụng \eqref{eq:ud1}, dựa vào phương trình trên ta sẽ thu được:
\begin{equation}\label{eq:ud4}
	dV^\tau(t)=[\lambda(t)V^\tau(t)+(\alpha(t)-\lambda(t))\tau_1(t)X_1(t)]dt +\beta(t)\tau_1(t)X_1(t)dB_t
\end{equation}
Đến đây ta thấy rằng phương trình \eqref{eq:ud4} là một phương trình vi phân ngẫu nhiên tuyến tính mà ta sẽ nói đến phương pháp giải ở mục phía sau.\\
Mặc khác, trong thực tế thường xét đến việc tìm phương án đầu tư $\tau(t);t\in[0,T]$ để thu được tổng giá trị tài sản là: $V^\tau(T)=\text{\AA}$. Đại lượng ngẫu nhiên $\text{\AA}$ tương ứng với quyền tài chính (financial claim), phương án đầu tư như vậy được gọi là phương án đáp ứng (replicating portfolio).\\
Mở rộng vấn đề nêu trên vào trường hợp nhiều chiều, tương ứng với việc có n khối tài sản rủi ro với biến động giá $X_i(t),i=1,2,...,n$, được cho bởi:
\begin{equation}\label{eq:ud5}
	dX_i(t)=X_i(t^-)\left[\beta_i(t)dB_t+\int_{R_0}\gamma_i(t,x)\overline{N}(dt,dx) \right] \condition{$i=1,2,...,n$}
\end{equation}
trong đó $B(t)=(B_1(t),...,B_n(t))^T;t\geq 0$ là quá trình Wiener n-chiều với:
\begin{equation*}
	\overline{N}(dt,dx)=\left(\overline{N}_1(dt,dx_1),...,\overline{N}_n(dt,dx_n)\right)^T \condition{$x\in R_0:=R\backslash\{0\}$}
\end{equation*}
là độ đo ngẫu nhiên Poisson bù (compensated Poisson random measure) và còn:
\begin{equation*}
	\beta_i(t)=(\beta_{i_1}(t),...,\beta_{i_n}(t));\gamma_i(t,x)=(\gamma_{i_1}(t,x),...,\gamma_{i_n}(t,x))
\end{equation*}
là các quá trình ngẫu nhiên thoả điều kiện:
\begin{equation*}
	\sum_{i,j=1}^{n}\int_{0}^{T}\left[\beta_{ij}^2(t)+\int_{R_0}\gamma_{ij}^2(t,x)v(dx) \right]dt<\infty \condition{h.c}
\end{equation*}
Quyền $\text{\AA}$ sẽ có thể đáp ứng nếu tồn tại quá trình ngẫu nhiên:
\begin{equation*}
	\varphi(t)=\left(\varphi_1(t),...\varphi_n(t) \right)^T \condition{$t\geq 0$ sao cho:}\\
\end{equation*}
\begin{equation*}
	\sum_{i=1}^{n}\int_{0}^{T}\varphi_i^2(t)X_i(t^-)\left[\sum_{j=1}^{n}\beta_{ij}^2(t)+\int_{R_0}\gamma_{ij}^2(t,x)v(dx) \right]dt<\infty \condition{h.c}
\end{equation*}
cùng với việc thoả điều kiện $\text{\AA}=E\{\text{\AA}\}+\sum\limits_{i=1}^{n}\int\limits_{0}^{T}\varphi_i(t)dX_i(t)$. Trong trường hợp này, ta nói rằng $\varphi$ là phương án đáp ứng đối với quyền tài chính.\\
\subsection{Giải quyết bài toán đầu tư đáp ứng}
Để giải quyết bài toán theo hướng đầu tư đáp ứng đã nêu, thường nó dẫn đến phương trình \eqref{eq:ud4} và \eqref{eq:ud5} mà ta xét chứng dưới dạng phương trình vi phân ngẫu nhiên tổng quát sau:\\
\textbf{Định lý:} Cho phương trình vi phân ngẫu nhiên tuyến tính có dạng:
\begin{dmath*}
dX(t)=[\alpha(t)X(t^-)+\mathscr{A}(t)]dt+[\beta(t)X(t^-)+\mathscr{B}(t)]dB_t\\ +\int_{R_0}[\gamma(t,x)+\mathscr{T}(t,x)]\overline{N}(dt,dx) \condition{$X(0)=x_0$}
\end{dmath*}
trong đó: $\alpha(T)=\alpha(t,\omega);\beta(t)=\beta(t,\omega);\gamma(t,x)=\gamma(t,x,\omega);\mathscr{A}(t)=\mathscr{A}(t,\omega);\mathscr{B}(t)=\mathscr{B}(t,\omega);\mathscr{T}(t,x)=\mathscr{T}(t,x,\omega)$ là những quá trình ngẫu nhiên thoả các điều kiện:
\begin{equation*}
\begin{split}
	&\int_{0}^{T}\left[\alpha(t)+\beta^2(t)+\int_{R_0}\gamma^2(t,x)v(dx) \right]dt<\infty\\
	&\int_{0}^{T}\left[\mathscr{A}(t)+\mathscr{B}^2(t)+\int_{R_0}\mathscr{T}^2(t,x)v(dx) \right]dt<\infty
\end{split}
\end{equation*}
Khi đó, ta sẽ có nghiệm của phương trình cho bởi hệ thức:
\begin{dmath*}
\dfrac{X(t)}{X_1(t^-)}=x_0+\int\limits_{0}^{t}X_1^{-1}(s^-)\left[\mathscr{A}(s)-\beta(s)\mathscr{B}(t)-\int\limits_{R_0}\dfrac{\gamma(s,x)\mathscr{T}(s,x)}{1+\gamma(s,x)}v(dx) \right]ds+\int\limits_{0}^{t}\dfrac{\mathscr{B}(s)}{X_1^{-1}(s^-)}dB(s)+\int\limits_{0}^{t}\int\limits_{R_0}\dfrac{\mathscr{T}(t,x)}{X_1(s^-)(1+\gamma(s,x))}\overline{N}(dt,dx).	
\end{dmath*}
trong đó:
\begin{dmath*}
X_1(t^-)=\exp\left\{\int_{0}^{t}\left[\alpha(s)-\dfrac{1}{2}\beta^2(t)+\int_{R_0}[\log(1+\gamma(s,x))-\gamma(s,x)]v(dx) \right]+\int_{0}^{t}\beta(s)dB(s)+\int_{R_0}\int_{0}^{t}\log(1+\gamma(s,x))\overline{N}(ds,dx) \right\}
\end{dmath*}
Để chứng minh định lý này ta dùng phương pháp tách nghiệm đã xét ở phần trước. Ngoài ra, để giải quyết cụ thể hơn các vấn đề về bài toán đầu tư rõ ràng ta phải xét thêm đến những loại quyền chọn cụ thể trên thị trường (quyền chọn Châu Âu, quyền chọn Châu Á...). Việc xem xét và giải chúng trong một vài trường hợp đã cho ta những kết quả nhất định như các kết quả kinh điển của Black-Scholes, Floyd B. Hanson.\\
\newpage
\section*{\begin{center}
	TÀI LIỆU THAM KHẢO
\end{center}}
\addcontentsline{toc}{section}{Tài liệu tham khảo}
\thispagestyle{empty}
\begin{enumerate}
	\item \textbf{Dương Tôn Đảm}, \textit{Quá trình ngẫu nhiên - Phần I: Tích phân và Phương trình vi phân ngẫu nhiên}, Nhà xuất bản Đại học Quốc gia Thành phố Hồ Chí Minh, Hồ Chí Minh, 2007.
	\item \textbf{Dương Tôn Đảm}, \textit{Quá trình ngẫu nhiên - Phần II: Các phép toán Malliavin}, Nhà xuất bản Đại học Quốc gia Thành phố Hồ Chí Minh, Hồ Chí Minh, 2010.
	\item \textbf{Dương Tôn Đảm}, \textit{Lớp quá trình ngẫu nhiên Itô - Levy và ứng dụng}, Nhà xuất bản Đại học Quốc gia Thành phố Hồ Chí Minh, Hồ Chí Minh, 2015.
	\item \textbf{Trần Hùng Thao}, \textit{Tích phân ngẫu nhiên \& phương trình vi phân ngẫu nhiên}, Nhà xuất bản Khoa học và Kỹ thuật, Hà Nội, 2000.
	\item \textbf{Bernt Oksendal}, \textit{Stochastic Differential Equations - An Introdution with Application}, Springer Verlag, 2006.
	\item \textbf{Krishna B. Athreya, Soumendra N. Lahiri}, Measure Theory and Probability Theory, Springer, New York, 2006.
\end{enumerate}
\newpage
\tableofcontents

%Nội dung trọng tâm
%\textbf{Phương trình vi phân ngẫu nhiên}\\
%Yêu cầu: Dạng phương trình vi phân ngẫu nhiên\\
%Nghiệm yếu và nghiệm mạnh\\
%Các dạng đặc biệt của phương trình vi phân ngẫu nhiên\\
%Phương trình vi phân tuyến tính.\\
\end{document}